%bismi ALLAH ARRAHMAN ARRAHIM













































% !TEX program = pdflatex
% ============================================================
% TH_SONDE — PROTOCOLE EXPERIMENTAL FONDAMENTAL
% CNG : Coeur Neuronal Général
% CNC : Coeur Neuronal de Calcul
%
% Source technique :
% sondes_ex7_v003.zip
% inputs/reference_v002/src/ex6_common.py
%
% Document autonome compilable avec pdfLaTeX.
% ============================================================

\documentclass[11pt,a4paper]{article}

% ------------------------------------------------------------
% LANGUE ET TYPOGRAPHIE
% ------------------------------------------------------------
\usepackage[utf8]{inputenc}
\usepackage[T1]{fontenc}
\usepackage[french]{babel}
\usepackage{lmodern}
\usepackage{microtype}

% ------------------------------------------------------------
% MISE EN PAGE
% ------------------------------------------------------------
\usepackage[
    a4paper,
    left=2cm,
    right=2cm,
    top=2cm,
    bottom=2cm
]{geometry}

\usepackage{setspace}
\setstretch{1.08}
\setlength{\parindent}{0pt}
\setlength{\parskip}{5pt}
\setlength{\emergencystretch}{3em}

% ------------------------------------------------------------
% MATHEMATIQUES
% ------------------------------------------------------------
\usepackage{amsmath,amssymb,amsthm,mathtools}
\usepackage{bm}

% ------------------------------------------------------------
% TABLEAUX
% ------------------------------------------------------------
\usepackage{booktabs}
\usepackage{tabularx}
\usepackage{array}
\usepackage{longtable}
\usepackage{makecell}
\usepackage{adjustbox}

\renewcommand{\arraystretch}{1.2}
\setlength{\tabcolsep}{4pt}

% ------------------------------------------------------------
% FIGURES ET GRAPHES
% ------------------------------------------------------------
\usepackage{graphicx}
\usepackage{float}
\usepackage{tikz}
\usetikzlibrary{
    arrows.meta,
    positioning,
    shapes.geometric,
    fit,
    calc
}

\usepackage{caption}
\captionsetup{
    font=small,
    labelfont=bf,
    width=0.95\textwidth
}

% ------------------------------------------------------------
% LISTES ET LIENS
% ------------------------------------------------------------
\usepackage{enumitem}
\setlist[itemize]{leftmargin=1.5em,itemsep=2pt}
\setlist[enumerate]{leftmargin=1.7em,itemsep=2pt}

\usepackage[hidelinks]{hyperref}

% ------------------------------------------------------------
% COMMANDES
% ------------------------------------------------------------
\newcommand{\R}{\mathbb{R}}
\newcommand{\N}{\mathbb{N}}
\newcommand{\Loss}{\mathcal{L}}
\newcommand{\CNG}{\textsc{CNG}}
\newcommand{\CNC}{\textsc{CNC}}
\newcommand{\RBF}{\textsc{RBF}}
\newcommand{\MSE}{\textsc{MSE}}
\newcommand{\MAE}{\textsc{MAE}}
\newcommand{\RMSE}{\textsc{RMSE}}

% ------------------------------------------------------------
% STYLES TIKZ
% ------------------------------------------------------------
\tikzset{
    block/.style={
        draw,
        rounded corners=2pt,
        align=center,
        minimum height=8mm,
        text width=25mm,
        inner sep=4pt,
        font=\small
    },
    smallblock/.style={
        draw,
        rounded corners=2pt,
        align=center,
        minimum height=7mm,
        text width=20mm,
        inner sep=3pt,
        font=\scriptsize
    },
    inputblock/.style={
        draw,
        rounded corners=2pt,
        align=center,
        minimum height=7mm,
        text width=16mm,
        inner sep=3pt,
        font=\small
    },
    flow/.style={
        -{Latex[length=2mm]},
        thick
    },
    optionalflow/.style={
        -{Latex[length=2mm]},
        dashed,
        thick
    }
}

% ============================================================
% DOCUMENT
% ============================================================

\title{
    \textbf{Fondations expérimentales de th\_sonde}\\
    \large Formalisation du coeur neuronal général et du coeur neuronal de calcul
}

\author{Projet th\_sonde}
\date{Version expérimentale 1.0}

\begin{document}

\maketitle

\begin{abstract}
Ce protocole définit deux expériences successives destinées
à tester les fondations de l'architecture th\_sonde.

La première expérience étudie un coeur neuronal général
(\CNG), composé de plusieurs sondes neuronales adaptatives
dont les sorties sont combinées pour approximer une fonction.

La deuxième expérience étudie un coeur neuronal de calcul
(\CNC), composé de sondes spécialisées dans des opérations
mathématiques élémentaires et capables d'être assemblées
dans des graphes de calcul.

La base technique est l'implémentation AdaptiveRBF
présente dans EX7 V003.

L'objectif n'est pas de démontrer immédiatement la
supériorité d'une architecture complète, mais de vérifier
séparément :
\begin{enumerate}
    \item la capacité d'une sonde à apprendre ;
    \item la capacité de plusieurs sondes à coopérer ;
    \item la capacité de sondes spécialisées à calculer ;
    \item la capacité de ces sondes à être composées ;
    \item le coût réel de cette composition.
\end{enumerate}

Aucun résultat expérimental nouveau n'est présupposé.
Les performances devront être mesurées selon le protocole
décrit dans ce document.
\end{abstract}

\tableofcontents
\newpage

% ============================================================
\section{Fondement technique : source EX7 V003}
% ============================================================

\subsection{Source de référence}

L'archive de référence est :

\begin{center}
\texttt{sondes\_ex7\_v003.zip}
\end{center}

Le fichier source principal utilisé est :

\begin{center}
\texttt{inputs/reference\_v002/src/ex6\_common.py}
\end{center}

La classe concernée est :

\begin{center}
\texttt{AdaptiveRBF}
\end{center}

Elle est définie aux lignes 115 à 138 du fichier.

\subsection{Structure réellement présente dans la source}

La sonde comprend :

\begin{enumerate}
    \item une projection linéaire apprise ;
    \item une activation hyperbolique tangente ;
    \item des centres RBF initialisés par KMeans ;
    \item des largeurs RBF apprenables ;
    \item une matrice de lecture apprenable ;
    \item un biais de sortie.
\end{enumerate}

L'encodeur de la source est :

\[
z(x)=\tanh(Wx+b)
\]

Les centres sont initialisés à partir de la projection
des données d'entraînement :

\[
C=\operatorname{KMeans}(z(X_{\mathrm{train}}))
\]

La distance au centre $k$ est :

\[
d_k^2(x)=\|z(x)-c_k\|_2^2
\]

Les largeurs sont paramétrées en logarithme :

\[
\sigma_k=\exp(\ell_k)
\]

La fonction d'activation RBF est :

\[
r_k(x)=
\exp\left(
-\frac{d_k^2(x)}{2\sigma_k^2}
\right)
\]

La sortie est :

\[
P(x)=r(x)^\top V+\beta
\]

Dans EX7, cette sortie possède deux dimensions,
correspondant aux deux logits de classification.

\subsection{Adaptation pour les expériences numériques}

Pour les expériences CNG et CNC, la dimension
de sortie doit être adaptée à la tâche.

Pour une fonction scalaire :

\[
n_{\mathrm{classes}}=1
\]

Pour une fonction vectorielle :

\[
n_{\mathrm{classes}}=q
\]

où $q$ est la dimension de sortie recherchée.

Il ne faut donc pas conserver automatiquement
les deux logits de la tâche Wisconsin.

La structure interne de la sonde est réutilisée,
mais la dimension de sa couche de lecture est adaptée.

\subsection{Différence entre la source et les nouvelles expériences}

\begin{tabularx}{\textwidth}{
    >{\RaggedRight\arraybackslash}p{0.32\textwidth}
    X
}
\toprule
\textbf{Élément} & \textbf{Statut} \\
\midrule
Encodeur Linear + tanh
& Présent dans la source. \\

Centres RBF KMeans
& Présents dans la source. \\

Largeurs apprenables
& Présentes dans la source. \\

Couche de lecture
& Présente dans la source. \\

Plusieurs sondes avec fusion apprise
& Nouvelle configuration CNG. \\

Sondes dédiées aux opérations
& Nouvelle configuration CNC. \\

Composition de graphes de calcul
& Nouvelle expérience CNC. \\

Coordinateur neuronal adaptatif
& Non inclus dans ces deux expériences. \\
\bottomrule
\end{tabularx}

\subsection{Précaution sur l'expérience EX7}

Dans EX7 V003 :

\begin{itemize}
    \item sept sondes sont entraînées ;
    \item les données supervisées sont réparties en sept shards ;
    \item les centres KMeans sont initialisés sur les caractéristiques
    de l'ensemble d'entraînement ;
    \item les sept sondes reçoivent les mêmes dimensions d'entrée ;
    \item les sept sondes sont exécutées à l'inférence ;
    \item le coordinateur $\rho$ est un vecteur global appris.
\end{itemize}

Il n'y a pas de sélection conditionnelle des sondes
à l'inférence.

Le présent protocole ne reprend pas automatiquement
le partitionnement des données en shards.

Pour le CNG initial, toutes les sondes reçoivent
le même jeu d'entraînement. Cette décision permet
d'isoler l'effet de la composition des sondes.

% ============================================================
\section{Question scientifique générale}
% ============================================================

\subsection{Question principale}

Peut-on construire un coeur fonctionnel à partir
de plusieurs sondes neuronales adaptatives, de telle sorte
que ces sondes apprennent collectivement une fonction
et puissent ensuite être spécialisées dans des opérations
réutilisables ?

\subsection{Décomposition de la question}

La question est décomposée en deux expériences.

\textbf{Expérience A : CNG}

Les sondes doivent apprendre une fonction générale
à partir d'exemples.

\textbf{Expérience B : CNC}

Les sondes doivent apprendre des opérations identifiables,
puis les composer dans un graphe de calcul.

Cette séparation évite de confondre :
\begin{itemize}
    \item capacité d'approximation ;
    \item coopération ;
    \item spécialisation ;
    \item composition ;
    \item coordination dynamique.
\end{itemize}

\subsection{Hypothèses}

\begin{tabularx}{\textwidth}{
    >{\centering\arraybackslash}p{0.10\textwidth}
    X
}
\toprule
\textbf{Code} & \textbf{Hypothèse falsifiable} \\
\midrule
H1
& Une sonde AdaptiveRBF peut approximer des fonctions
numériques non linéaires sur un domaine défini. \\

H2
& Plusieurs sondes combinées peuvent apprendre
une fonction avec une erreur de généralisation contrôlée. \\

H3
& La composition apprise de plusieurs sondes peut
apporter un bénéfice mesurable par rapport à une sonde unique. \\

H4
& Une sonde spécialisée peut apprendre une opération
mathématique élémentaire. \\

H5
& Des sondes opérationnelles entraînées séparément
peuvent être composées dans un graphe de calcul inédit. \\

H6
& Le coût de composition reste mesurable et ne doit pas
annuler le bénéfice fonctionnel obtenu. \\
\bottomrule
\end{tabularx}

Aucune de ces hypothèses ne doit être considérée
comme démontrée avant l'exécution des expériences.

% ============================================================
\section{Définition mathématique d'une sonde}
% ============================================================

\subsection{Entrée et paramètres}

Une sonde $P_i$ reçoit :

\[
x\in\R^n
\]

Ses paramètres sont :

\[
\theta_i=
(W_i,b_i,C_i,\ell_i,V_i,\beta_i)
\]

avec :

\[
W_i\in\R^{d\times n}
\]

\[
b_i\in\R^d
\]

\[
C_i\in\R^{K\times d}
\]

\[
\ell_i\in\R^K
\]

\[
V_i\in\R^{K\times q}
\]

\[
\beta_i\in\R^q
\]

où :
\begin{itemize}
    \item $n$ est la dimension de l'entrée ;
    \item $d$ est la dimension latente ;
    \item $K$ est le nombre de centres RBF ;
    \item $q$ est la dimension de sortie.
\end{itemize}

\subsection{Encodeur latent}

\[
z_i(x)=\tanh(W_i x+b_i)
\]

\subsection{Activation RBF}

\[
\sigma_{i,k}=\exp(\ell_{i,k})
\]

\[
r_{i,k}(x)=
\exp\left(
-\frac{
\|z_i(x)-c_{i,k}\|_2^2
}{
2\sigma_{i,k}^2
}
\right)
\]

Les largeurs doivent respecter les bornes numériques
définies dans l'implémentation :

\[
0.05\leq\sigma_{i,k}\leq12
\]

\subsection{Lecture}

\[
P_i(x;\theta_i)
=
r_i(x)^\top V_i+\beta_i
\]

\subsection{Initialisation}

Les centres sont initialisés avec KMeans sur les
représentations latentes du jeu d'entraînement uniquement.

Les données de validation et de test sont exclues
de cette initialisation.

Pour chaque sonde :
\begin{enumerate}
    \item initialiser $W_i$ et $b_i$ ;
    \item calculer $z_i(X_{\mathrm{train}})$ ;
    \item appliquer KMeans ;
    \item initialiser les centres $C_i$ ;
    \item calculer une largeur initiale à partir
    des distances inter-centres ;
    \item initialiser $V_i$ et $\beta_i$ ;
    \item entraîner les paramètres autorisés.
\end{enumerate}

La méthode exacte d'initialisation est enregistrée
dans la configuration de chaque expérience.

\subsection{Graphe interne d'une sonde}

\begin{figure}[H]
\centering
\begin{adjustbox}{max width=\textwidth}
\begin{tikzpicture}[node distance=7mm and 9mm]

\node[inputblock] (x) {Entrée $x$};

\node[block, right=of x] (encoder)
{Encodeur latent\\
$z=\tanh(Wx+b)$};

\node[block, right=of encoder] (rbf)
{Activations RBF\\
$r_1,\ldots,r_K$};

\node[block, right=of rbf] (readout)
{Lecture\\
$r^\top V+\beta$};

\node[inputblock, right=of readout] (y)
{Sortie $P_i(x)$};

\draw[flow] (x)--(encoder);
\draw[flow] (encoder)--(rbf);
\draw[flow] (rbf)--(readout);
\draw[flow] (readout)--(y);

\end{tikzpicture}
\end{adjustbox}
\caption{Structure interne d'une sonde adaptative dérivée d'EX7.}
\end{figure}

% ============================================================
\section{Expérience A : coeur neuronal général}
% ============================================================

\subsection{Définition du CNG}

Le coeur neuronal général est un ensemble de $M$ sondes :

\[
\mathcal{C}_{G}
=
\{P_1,P_2,\ldots,P_M\}
\]

Chaque sonde reçoit la même entrée :

\[
h_i=P_i(x;\theta_i)
\]

Les sorties sont concaténées :

\[
h(x)=
\begin{bmatrix}
h_1(x)\\
h_2(x)\\
\vdots\\
h_M(x)
\end{bmatrix}
\]

Un composeur apprend à combiner ces sorties.

Pour la première expérience, le composeur est linéaire :

\[
C_\phi(h)=W_c h+b_c
\]

La fonction globale du CNG est :

\[
\boxed{
\widehat y(x)=
W_c
\begin{bmatrix}
P_1(x;\theta_1)\\
P_2(x;\theta_2)\\
\vdots\\
P_M(x;\theta_M)
\end{bmatrix}
+b_c
}
\]

Les paramètres optimisés sont :

\[
\Theta=
\{\theta_1,\ldots,\theta_M,W_c,b_c\}
\]

\subsection{Architecture du CNG}

\begin{figure}[H]
\centering
\begin{adjustbox}{max width=\textwidth}
\begin{tikzpicture}[node distance=8mm and 9mm]

\node[inputblock] (x) {Entrée $x$};

\node[smallblock, right=12mm of x, yshift=16mm] (p1)
{Sonde $P_1$};

\node[smallblock, right=12mm of x] (p2)
{Sonde $P_2$};

\node[smallblock, right=12mm of x, yshift=-16mm] (pm)
{Sonde $P_M$};

\node[block, right=15mm of p2] (concat)
{Concaténation\\des sorties};

\node[block, right=of concat] (compose)
{Composeur appris\\$C_\phi$};

\node[inputblock, right=of compose] (y)
{Sortie $\widehat y$};

\draw[flow] (x.east)--(p1.west);
\draw[flow] (x.east)--(p2.west);
\draw[flow] (x.east)--(pm.west);

\draw[flow] (p1.east)--(concat.west);
\draw[flow] (p2.east)--(concat.west);
\draw[flow] (pm.east)--(concat.west);

\draw[flow] (concat)--(compose);
\draw[flow] (compose)--(y);

\end{tikzpicture}
\end{adjustbox}
\caption{Graphe du coeur neuronal général.}
\end{figure}

\subsection{Ce qui est volontairement absent}

La version CNG-1 ne contient pas :
\begin{itemize}
    \item de coordinateur externe ;
    \item de sélection dynamique de sondes ;
    \item de routage conditionnel ;
    \item de mémoire temporaire ;
    \item de modification du graphe ;
    \item de communication récurrente entre sondes.
\end{itemize}

La seule coopération étudiée est la combinaison
apprise de leurs sorties.

Cette restriction permet d'isoler l'effet du
multi-sondes.

\subsection{Fonctions à apprendre}

Quatre fonctions synthétiques sont retenues.

\[
f_1(x)=\sin(x)
\]

\[
f_2(x_1,x_2)=x_1x_2
\]

\[
f_3(x_1,x_2)=x_1^2+x_2
\]

\[
f_4(x_1,x_2,x_3)
=
\sin(x_1x_2)+x_3
\]

Elles couvrent :

\begin{tabularx}{\textwidth}{
    >{\centering\arraybackslash}p{0.13\textwidth}
    X
}
\toprule
\textbf{Fonction} & \textbf{Propriété testée} \\
\midrule
$f_1$ & Approximation non linéaire simple. \\
$f_2$ & Interaction multiplicative. \\
$f_3$ & Composition quadratique et additive. \\
$f_4$ & Interaction non linéaire de plusieurs variables. \\
\bottomrule
\end{tabularx}

\subsection{Domaines}

Les domaines d'entraînement et d'interpolation sont :

\[
f_1:\quad x\in[-\pi,\pi]
\]

\[
f_2:\quad (x_1,x_2)\in[-2,2]^2
\]

\[
f_3:\quad (x_1,x_2)\in[-2,2]^2
\]

\[
f_4:\quad (x_1,x_2,x_3)\in[-2,2]^3
\]

Un jeu d'extrapolation séparé est généré sur :

\[
[-3,3]^n
\]

Les résultats d'interpolation et d'extrapolation
ne doivent jamais être fusionnés dans une seule métrique.

\subsection{Génération et séparation des données}

Pour chaque fonction :

\[
D=
D_{\mathrm{train}}
\cup
D_{\mathrm{val}}
\cup
D_{\mathrm{test}}
\]

Les ensembles sont disjoints.

Répartition proposée :

\begin{center}
\begin{tabular}{lc}
\toprule
Ensemble & Proportion \\
\midrule
Entraînement & 70 \% \\
Validation & 15 \% \\
Test & 15 \% \\
\bottomrule
\end{tabular}
\end{center}

Le jeu de test est généré indépendamment du jeu
d'entraînement et du jeu de validation.

Le test ne doit pas être consulté pendant
la sélection des hyperparamètres.

Les données de validation et de test ne sont pas
utilisées pour initialiser les centres KMeans.

\subsection{Configurations à comparer}

\begin{tabularx}{\textwidth}{
    >{\centering\arraybackslash}p{0.10\textwidth}
    >{\RaggedRight\arraybackslash}p{0.34\textwidth}
    X
}
\toprule
\textbf{ID} & \textbf{Configuration} & \textbf{Objectif} \\
\midrule
G1
& Une sonde AdaptiveRBF
& Référence individuelle. \\

G2
& Plusieurs sondes avec fusion moyenne fixe
& Effet du nombre de sondes sans apprentissage du composeur. \\

G3
& Plusieurs sondes avec composeur appris
& Architecture CNG candidate. \\

G4
& MLP dense de capacité comparable
& Comparateur neuronal classique. \\
\bottomrule
\end{tabularx}

La fusion fixe de G2 est :

\[
\widehat y_{\mathrm{fixe}}
=
\frac{1}{M}
\sum_{i=1}^{M}P_i(x)
\]

La fusion apprise de G3 est :

\[
\widehat y_{\mathrm{apprise}}
=
W_c[P_1(x);\ldots;P_M(x)]+b_c
\]

\subsection{Tailles initiales}

La configuration principale reprend les dimensions
de référence de la sonde EX7 :

\[
d=4,\qquad K=8
\]

Le nombre de sondes est :

\[
M\in\{1,2,4,7\}
\]

Les expériences secondaires peuvent étudier :

\[
d\in\{2,4,8\}
\]

\[
K\in\{4,8,16\}
\]

La recherche secondaire ne doit commencer qu'après
la réalisation de l'expérience principale.

\subsection{Fonction de perte}

Pour une sortie scalaire :

\[
\Loss_{\mathrm{CNG}}
=
\frac{1}{N}
\sum_{j=1}^{N}
\left(
\widehat y_j-y_j
\right)^2
\]

Pour une sortie de dimension $q$ :

\[
\Loss_{\mathrm{CNG}}
=
\frac{1}{Nq}
\sum_{j=1}^{N}
\sum_{k=1}^{q}
(\widehat y_{j,k}-y_{j,k})^2
\]

L'ensemble des sondes et le composeur sont entraînés
conjointement.

\subsection{Optimisation}

Configuration initiale :

\begin{itemize}
    \item Optimiseur : Adam.
    \item Learning rate : $10^{-3}$.
    \item Weight decay : $10^{-5}$.
    \item Batch size : 64.
    \item Maximum : 500 époques.
    \item Early stopping : patience 30.
    \item Sélection : MSE de validation.
\end{itemize}

Ces valeurs constituent une configuration initiale
à geler avant l'expérience principale.

Les résultats doivent également contenir les valeurs
réellement utilisées par le programme.

\subsection{Mesure de la coopération}

Le nombre de sondes ne prouve pas leur coopération.

Pour mesurer la contribution d'une sonde $i$,
on compare le modèle complet au modèle dont cette
sonde a été retirée.

Modèle complet :

\[
\widehat y_{\mathrm{full}}(x)
=
C_\phi(P_1(x),\ldots,P_M(x))
\]

Modèle sans sonde $i$ :

\[
\widehat y_{-i}(x)
=
C_{\phi,-i}(P_1(x),\ldots,P_{i-1}(x),
P_{i+1}(x),\ldots,P_M(x))
\]

Le composeur réduit est réentraîné dans une expérience
d'ablation indépendante.

On calcule :

\[
\Delta_i=
\Loss_{-i}-\Loss_{\mathrm{full}}
\]

Une valeur positive signifie que le retrait de la sonde
dégrade la performance.

Une valeur nulle ou négative indique que la sonde
n'apporte pas de contribution positive mesurable
dans cette configuration.

On complète cette analyse par :
\begin{itemize}
    \item la corrélation des sorties des sondes ;
    \item la corrélation de leurs erreurs ;
    \item la contribution de chaque sonde ;
    \item la stabilité de ces mesures entre graines.
\end{itemize}

Il ne faut pas conclure à une spécialisation uniquement
parce que les sondes ont des paramètres indépendants.

% ============================================================
\section{Expérience B : coeur neuronal de calcul}
% ============================================================

\subsection{Définition}

Le coeur neuronal de calcul est un ensemble de sondes
dont chacune apprend une opération mathématique.

On définit :

\[
\mathcal{O}
=
\{
\mathrm{ADD},
\mathrm{SUB},
\mathrm{MUL},
\mathrm{DIV},
\mathrm{EXP},
\mathrm{LOG}
\}
\]

Chaque opération possède une sonde :

\[
P_o(u;\theta_o)
\]

La sonde doit approximer la fonction mathématique
associée à l'opération.

Exemples :

\[
P_{\mathrm{ADD}}(a,b)\approx a+b
\]

\[
P_{\mathrm{MUL}}(a,b)\approx ab
\]

\[
P_{\mathrm{EXP}}(a)\approx e^a
\]

\[
P_{\mathrm{LOG}}(a)\approx\log(a)
\]

\subsection{Domaines de validité}

Les domaines suivants sont fixés pour l'expérience
principale.

\begin{tabularx}{\textwidth}{
    >{\centering\arraybackslash}p{0.16\textwidth}
    X
}
\toprule
\textbf{Opération} & \textbf{Domaine} \\
\midrule
ADD & $(a,b)\in[-2,2]^2$. \\

SUB & $(a,b)\in[-2,2]^2$. \\

MUL & $(a,b)\in[-2,2]^2$. \\

DIV & $(a,b)\in[-2,2]^2$ avec $|b|\geq0.25$. \\

EXP & $a\in[-2,2]$. \\

LOG & $a\in[0.1,4]$. \\
\bottomrule
\end{tabularx}

Ces restrictions sont obligatoires.

Une division par zéro n'est pas une entrée valide.

Un logarithme d'un nombre négatif n'est pas une entrée valide.

Les expériences hors domaine sont rapportées comme
des tests de robustesse distincts, jamais comme des
évaluations normales.

\subsection{Apprentissage individuel des opérations}

Pour chaque opération $o$, on génère :

\[
D_o=
\{(u_j,o(u_j))\}_{j=1}^{N_o}
\]

La perte de la sonde est :

\[
\Loss_o(\theta_o)
=
\frac{1}{N_o}
\sum_{j=1}^{N_o}
\left(
P_o(u_j;\theta_o)-o(u_j)
\right)^2
\]

Chaque sonde est entraînée séparément.

Les données sont divisées en trois ensembles :
\[
D_o=
D_{o,\mathrm{train}}
\cup D_{o,\mathrm{val}}
\cup D_{o,\mathrm{test}}
\]

Les centres KMeans sont initialisés uniquement
sur les représentations des données d'entraînement.

\subsection{Architecture du CNC}

\begin{figure}[H]
\centering
\begin{adjustbox}{max width=\textwidth}
\begin{tikzpicture}[node distance=8mm and 9mm]

\node[block] (input) {Entrées\\numériques};

\node[smallblock, right=13mm of input, yshift=22mm] (add)
{Sonde ADD};

\node[smallblock, right=13mm of input, yshift=11mm] (sub)
{Sonde SUB};

\node[smallblock, right=13mm of input] (mul)
{Sonde MUL};

\node[smallblock, right=13mm of input, yshift=-11mm] (div)
{Sonde DIV};

\node[smallblock, right=13mm of input, yshift=-22mm] (exp)
{Sonde EXP / LOG};

\node[block, right=17mm of mul] (graph)
{Graphe de calcul\\explicite};

\node[block, right=of graph] (result)
{Résultat numérique};

\draw[flow] (input.east)--(add.west);
\draw[flow] (input.east)--(sub.west);
\draw[flow] (input.east)--(mul.west);
\draw[flow] (input.east)--(div.west);
\draw[flow] (input.east)--(exp.west);

\draw[flow] (add.east)--(graph.west);
\draw[flow] (sub.east)--(graph.west);
\draw[flow] (mul.east)--(graph.west);
\draw[flow] (div.east)--(graph.west);
\draw[flow] (exp.east)--(graph.west);

\draw[flow] (graph)--(result);

\end{tikzpicture}
\end{adjustbox}
\caption{Organisation des sondes opérationnelles du CNC.}
\end{figure}

Le graphe de calcul est explicite dans cette première
version.

Le CNC ne décide pas encore de l'opération à exécuter.

Cette décision est volontaire : elle permet de tester
la capacité des sondes à calculer avant de tester
la capacité d'un coordinateur à choisir les opérations.

\subsection{Premier graphe composé}

La première expression composée est :

\[
f(a,b,c)=ab+e^c
\]

Les opérations sont :

\[
u_1=P_{\mathrm{MUL}}(a,b)
\]

\[
u_2=P_{\mathrm{EXP}}(c)
\]

\[
\widehat y=P_{\mathrm{ADD}}(u_1,u_2)
\]

La cible exacte est :

\[
y=ab+e^c
\]

\subsection{Graphe du premier calcul}

\begin{figure}[H]
\centering
\begin{adjustbox}{max width=\textwidth}
\begin{tikzpicture}[node distance=8mm and 12mm]

\node[inputblock] (a) {$a$};
\node[inputblock, below=of a] (b) {$b$};
\node[inputblock, below=of b] (c) {$c$};

\node[block, right=of a, yshift=-4mm] (mul)
{Sonde MUL\\$u_1=P_{\mathrm{MUL}}(a,b)$};

\node[block, right=of c] (exp)
{Sonde EXP\\$u_2=P_{\mathrm{EXP}}(c)$};

\node[block, right=19mm of mul, yshift=-9mm] (add)
{Sonde ADD\\$\widehat y=P_{\mathrm{ADD}}(u_1,u_2)$};

\node[inputblock, right=of add] (out)
{$\widehat y$};

\draw[flow] (a)--(mul);
\draw[flow] (b)--(mul);
\draw[flow] (c)--(exp);
\draw[flow] (mul)--(add);
\draw[flow] (exp.east)--([yshift=-5mm]add.west);
\draw[flow] (add)--(out);

\end{tikzpicture}
\end{adjustbox}
\caption{Graphe neuronal du calcul $ab+e^c$.}
\end{figure}

\subsection{Deuxième niveau de composition}

Après validation du premier graphe, trois expressions
sont retenues :

\[
f_1(a,b,c)=ab+e^c
\]

\[
f_2(a,b,c)=\frac{ab}{c+1}
\]

\[
f_3(a,b,c)=
\log(a^2+b^2+1)+e^c
\]

Les graphes doivent être définis avant l'évaluation.

Pour $f_2$, on impose :

\[
|c+1|\geq0.25
\]

Pour $f_3$, l'argument du logarithme est strictement positif.

\subsection{Propagation des erreurs}

Chaque opération apprise possède une erreur locale :

\[
\epsilon_o(u)=P_o(u)-o(u)
\]

L'erreur finale est :

\[
\epsilon_{\mathrm{final}}
=
\widehat y-y
\]

Pour le premier graphe :

\[
\widehat y=
P_{\mathrm{ADD}}
\left(
P_{\mathrm{MUL}}(a,b),
P_{\mathrm{EXP}}(c)
\right)
\]

L'erreur finale dépend :
\begin{itemize}
    \item de l'erreur de multiplication ;
    \item de l'erreur exponentielle ;
    \item de l'erreur d'addition ;
    \item de la sensibilité des opérations suivantes
    aux erreurs des opérations précédentes.
\end{itemize}

Il faut donc enregistrer les sorties intermédiaires
et calculer les erreurs à chaque noeud.

Une bonne précision sur chaque opération isolée
ne garantit pas une bonne précision sur un graphe composé.

\subsection{Trois configurations CNC}

\begin{tabularx}{\textwidth}{
    >{\centering\arraybackslash}p{0.10\textwidth}
    >{\RaggedRight\arraybackslash}p{0.33\textwidth}
    X
}
\toprule
\textbf{ID} & \textbf{Configuration} & \textbf{Objectif} \\
\midrule
C1
& Sondes opérationnelles entraînées séparément,
puis gelées
& Mesurer la composition sans ajustement global. \\

C2
& Sondes opérationnelles avec ajustement conjoint
sur les graphes
& Mesurer l'effet d'un ajustement global. \\

C3
& Modèle neuronal direct entraîné sur chaque expression
& Comparateur sans composition opérationnelle explicite. \\
\bottomrule
\end{tabularx}

C1 est l'expérience principale.

C2 est secondaire.

C3 constitue le comparateur direct.

\subsection{Pourquoi C1 est prioritaire}

Si les sondes opérationnelles sont entraînées séparément,
puis assemblées sans réentraînement, une réussite
sur des graphes non vus constitue une preuve plus forte
de réutilisabilité.

Si l'on entraîne directement l'ensemble du graphe,
il devient difficile de savoir si le modèle a appris
des opérations réutilisables ou simplement une fonction
globale.

C'est pourquoi C1 doit être exécutée avant C2.

% ============================================================
\section{Comparateurs et équité expérimentale}
% ============================================================

\subsection{Comparateurs CNG}

Les comparateurs sont :

\begin{enumerate}
    \item une sonde AdaptiveRBF unique ;
    \item plusieurs sondes à fusion fixe ;
    \item plusieurs sondes à fusion apprise ;
    \item un MLP de capacité comparable.
\end{enumerate}

\subsection{Comparateurs CNC}

Les comparateurs sont :

\begin{enumerate}
    \item calcul natif exact ;
    \item MLP direct sur l'expression ;
    \item sonde AdaptiveRBF directe sur l'expression ;
    \item CNC composé avec sondes gelées ;
    \item CNC composé avec ajustement conjoint.
\end{enumerate}

Le calcul natif est une référence exacte et une référence
de coût. Il ne constitue pas un concurrent neuronal
de capacité comparable.

\subsection{Contrôle du nombre de paramètres}

Pour chaque modèle :

\[
P_{\mathrm{total}}
=
P_{\mathrm{entraînable}}
+
P_{\mathrm{buffers}}
+
P_{\mathrm{prétraitement}}
\]

Les paramètres doivent être comptés automatiquement.

Le rapport doit donner :
\begin{itemize}
    \item le nombre de paramètres entraînables ;
    \item le nombre de paramètres fixes ;
    \item le nombre de paramètres de prétraitement ;
    \item le total persistant ;
    \item la taille du checkpoint.
\end{itemize}

Les comparaisons de qualité doivent être accompagnées
du nombre de paramètres.

Une différence de capacité importante doit être
explicitement signalée.

\subsection{Graines}

On reprend les cinq graines utilisées dans EX7 :

\[
S=\{1501,1607,1717,2029,3037\}
\]

Elles permettent une première analyse de sensibilité
à l'initialisation.

Elles ne constituent pas cinq jeux de données indépendants.

Pour une validation confirmatoire, il faudra ultérieurement
augmenter le nombre de répétitions et varier les
échantillons de données.

% ============================================================
\section{Mesures et métriques}
% ============================================================

\subsection{Qualité numérique}

Erreur absolue moyenne :

\[
\MAE=
\frac{1}{N}
\sum_{i=1}^{N}
|\widehat y_i-y_i|
\]

Erreur quadratique moyenne :

\[
\MSE=
\frac{1}{N}
\sum_{i=1}^{N}
(\widehat y_i-y_i)^2
\]

Racine de l'erreur quadratique moyenne :

\[
\RMSE=\sqrt{\MSE}
\]

Erreur relative moyenne :

\[
E_{\mathrm{rel}}
=
\frac{1}{N}
\sum_{i=1}^{N}
\frac{|\widehat y_i-y_i|}
{\max(|y_i|,\varepsilon)}
\]

La valeur de $\varepsilon$ doit être fixée avant
l'évaluation.

L'erreur relative ne doit pas être utilisée seule
lorsque les cibles sont proches de zéro.

\subsection{Coût d'exécution}

Pour chaque architecture, mesurer :

\begin{itemize}
    \item latence p50 ;
    \item latence p95 ;
    \item latence p99 ;
    \item débit ;
    \item nombre de paramètres ;
    \item mémoire persistante ;
    \item mémoire maximale, si mesurable ;
    \item durée d'entraînement.
\end{itemize}

Les latences sont mesurées dans des conditions identiques :
\begin{itemize}
    \item même matériel ;
    \item même nombre de threads ;
    \item même précision numérique ;
    \item même taille de batch ;
    \item même nombre de warmups ;
    \item même nombre de répétitions.
\end{itemize}

Les temps de construction et d'initialisation
sont rapportés séparément du temps d'inférence.

\subsection{Mesures propres au CNC}

Pour chaque graphe :

\[
N_{\mathrm{sondes}}
=
\text{nombre de sondes appelées}
\]

\[
N_{\mathrm{noeuds}}
=
\text{nombre de noeuds du graphe}
\]

\[
D_{\mathrm{graphe}}
=
\text{profondeur maximale}
\]

Il faut aussi mesurer :

\[
T_{\mathrm{total}}
=
T_{\mathrm{sondes}}
+
T_{\mathrm{composition}}
\]

Cette décomposition permet de vérifier si la composition
introduit un coût disproportionné.

% ============================================================
\section{Plan expérimental complet}
% ============================================================

\subsection{Séquence des expériences}

\begin{figure}[H]
\centering
\begin{adjustbox}{max width=0.78\textwidth}
\begin{tikzpicture}[node distance=6mm]

\node[block] (a)
{Gel du protocole};

\node[block, below=of a] (b)
{Validation logicielle};

\node[block, below=of b] (c)
{Entraînement CNG};

\node[block, below=of c] (d)
{Ablations CNG};

\node[block, below=of d] (e)
{Entraînement des opérations CNC};

\node[block, below=of e] (f)
{Composition des graphes};

\node[block, below=of f] (g)
{Tests de généralisation};

\node[block, below=of g] (h)
{Audit et décision};

\draw[flow] (a)--(b);
\draw[flow] (b)--(c);
\draw[flow] (c)--(d);
\draw[flow] (d)--(e);
\draw[flow] (e)--(f);
\draw[flow] (f)--(g);
\draw[flow] (g)--(h);

\end{tikzpicture}
\end{adjustbox}
\caption{Ordre obligatoire des opérations expérimentales.}
\end{figure}

\subsection{Phase 0 : gel du protocole}

Avant l'entraînement :
\begin{enumerate}
    \item fixer les fonctions ;
    \item fixer les domaines ;
    \item fixer les graines ;
    \item fixer les architectures ;
    \item fixer les hyperparamètres ;
    \item fixer les métriques ;
    \item enregistrer les configurations.
\end{enumerate}

Toute modification après consultation des résultats
de test doit produire une nouvelle version du protocole.

\subsection{Phase 1 : validation logicielle}

Les tests unitaires doivent vérifier :

\begin{itemize}
    \item les dimensions des tenseurs ;
    \item la stabilité des activations RBF ;
    \item les bornes des largeurs ;
    \item la reproductibilité des initialisations ;
    \item l'absence de valeurs NaN ;
    \item l'absence de valeurs infinies ;
    \item la séparation des ensembles ;
    \item le calcul exact des opérations de référence.
\end{itemize}

\subsection{Phase 2 : exécution CNG}

Pour chaque fonction et chaque graine :

\begin{enumerate}
    \item générer les données ;
    \item séparer train, validation et test ;
    \item initialiser les sondes sur le train ;
    \item initialiser le composeur ;
    \item entraîner le modèle ;
    \item sélectionner le checkpoint sur validation ;
    \item évaluer sur test ;
    \item enregistrer prédictions et métriques.
\end{enumerate}

\subsection{Phase 3 : ablations CNG}

Pour chaque sonde :
\begin{enumerate}
    \item retirer la sonde ;
    \item reconstruire le composeur ;
    \item réentraîner le modèle réduit ;
    \item comparer les métriques ;
    \item enregistrer la contribution.
\end{enumerate}

Les résultats d'ablation doivent être comparés
au modèle complet sur les mêmes partitions.

\subsection{Phase 4 : entraînement des opérations CNC}

Pour chaque opération :

\begin{enumerate}
    \item générer les entrées ;
    \item calculer les sorties exactes ;
    \item séparer les ensembles ;
    \item initialiser la sonde ;
    \item entraîner la sonde ;
    \item sélectionner le checkpoint ;
    \item évaluer l'opération ;
    \item sauvegarder le modèle.
\end{enumerate}

\subsection{Phase 5 : composition CNC}

Les sondes sont chargées sans modification.

Pour chaque graphe :
\begin{enumerate}
    \item fournir les entrées ;
    \item exécuter les opérations dans l'ordre ;
    \item sauvegarder les sorties intermédiaires ;
    \item comparer chaque noeud à sa référence exacte ;
    \item comparer la sortie finale à la cible ;
    \item mesurer le coût d'exécution.
\end{enumerate}

\subsection{Phase 6 : généralisation}

Le CNC est testé sur des compositions qui n'ont pas
été utilisées pendant l'entraînement.

Les tests comprennent :
\begin{itemize}
    \item des entrées numériques inédites ;
    \item des compositions inédites ;
    \item des valeurs proches des limites du domaine ;
    \item des domaines d'extrapolation séparés.
\end{itemize}

\subsection{Phase 7 : audit final}

Les éléments suivants sont conservés :

\begin{itemize}
    \item code source ;
    \item configuration ;
    \item graines ;
    \item checkpoints ;
    \item prédictions ;
    \item métriques par graine ;
    \item historiques d'entraînement ;
    \item profils de latence ;
    \item journaux d'erreurs ;
    \item empreintes SHA-256.
\end{itemize}

Les expériences échouées ne sont pas supprimées.

% ============================================================
\section{Critères de validation}
% ============================================================

\subsection{Niveau 1 : fonctionnement individuel}

Une sonde est fonctionnelle si :
\begin{enumerate}
    \item elle apprend l'opération ou la fonction ciblée ;
    \item elle converge de manière reproductible ;
    \item elle généralise à des entrées non vues ;
    \item elle reste numériquement stable.
\end{enumerate}

\subsection{Niveau 2 : fonctionnement collectif}

Le CNG est fonctionnel si le modèle multi-sondes
apprend correctement les fonctions testées.

La coopération est validée seulement si elle apporte
un bénéfice mesurable par rapport à une sonde unique
et aux comparateurs à capacité contrôlée.

\subsection{Niveau 3 : composition}

Le CNC est fonctionnel si :
\begin{enumerate}
    \item les opérations individuelles sont apprises ;
    \item les opérations peuvent être assemblées ;
    \item les graphes composés produisent les sorties attendues ;
    \item les erreurs intermédiaires sont maîtrisées ;
    \item les compositions inédites sont évaluées.
\end{enumerate}

\subsection{Niveau 4 : efficacité}

Une bonne précision ne prouve pas une meilleure efficacité.

Une revendication d'efficacité exige :
\begin{itemize}
    \item une mesure de latence ;
    \item une mesure du nombre d'opérations ;
    \item une mesure de mémoire ;
    \item une comparaison avec les références ;
    \item une qualité numérique suffisante.
\end{itemize}

Le calcul natif exact est une référence particulièrement
importante pour le CNC.

Le CNC ne sera pas déclaré plus rapide qu'un calcul natif
sans mesure directe.

\subsection{Tableau de décision}

\begin{tabularx}{\textwidth}{
    >{\RaggedRight\arraybackslash}p{0.34\textwidth}
    X
}
\toprule
\textbf{Résultat} & \textbf{Décision} \\
\midrule
La sonde individuelle échoue
& Corriger la sonde avant toute composition. \\

La sonde fonctionne mais le CNG échoue
& Étudier le composeur, l'initialisation et l'interaction
entre les sondes. \\

Le CNG fonctionne sans bénéfice multi-sondes
& Conserver le résultat comme approximation distribuée,
sans revendiquer un avantage collectif. \\

Le CNG montre un bénéfice reproductible
& Autoriser l'expérience CNC. \\

Les opérations CNC fonctionnent individuellement,
mais les graphes échouent
& Étudier la propagation des erreurs et la stabilité. \\

Les graphes fonctionnent sans réentraînement
& Valider la réutilisation des sondes opérationnelles
dans le périmètre testé. \\

La qualité est bonne mais le coût est supérieur
& Conclure à une capacité fonctionnelle, pas à une
amélioration d'efficacité. \\
\bottomrule
\end{tabularx}

% ============================================================
\section{Analyse statistique}
% ============================================================

\subsection{Résultats par graine}

Pour une métrique $m$ et $S$ graines :

\[
\overline m=
\frac{1}{S}\sum_{s=1}^{S}m_s
\]

Écart-type :

\[
s_m=
\sqrt{
\frac{1}{S-1}
\sum_{s=1}^{S}(m_s-\overline m)^2
}
\]

Les résultats sont présentés sous la forme :

\[
\overline m\pm s_m
\]

Les résultats individuels restent accessibles.

\subsection{Comparaison appariée}

Pour deux modèles A et B :

\[
\Delta_s=m_{A,s}-m_{B,s}
\]

\[
\overline\Delta=
\frac{1}{S}\sum_{s=1}^{S}\Delta_s
\]

Les graines communes permettent une comparaison
appariée de l'initialisation.

Elles ne remplacent pas des échantillons de données
indépendants.

Les intervalles de confiance obtenus avec cinq graines
doivent être interprétés comme descriptifs et non
comme une preuve confirmatoire.

% ============================================================
\section{Organisation du code}
% ============================================================

\subsection{Structure du projet}

\begin{verbatim}
th_sonde/
|
+-- configs/
|   +-- cng_config.json
|   +-- cnc_config.json
|   +-- seeds.json
|
+-- src/
|   +-- adaptive_rbf.py
|   +-- neural_core_general.py
|   +-- neural_core_calculation.py
|   +-- synthetic_data.py
|   +-- operation_graph.py
|   +-- train_cng.py
|   +-- train_cnc_operations.py
|   +-- evaluate_cnc_graphs.py
|   +-- benchmark.py
|   +-- audit.py
|
+-- tests/
|   +-- test_adaptive_rbf.py
|   +-- test_cng.py
|   +-- test_operations.py
|   +-- test_graphs.py
|   +-- test_data_integrity.py
|
+-- results/
|   +-- raw/
|   +-- models/
|   +-- metrics/
|   +-- profiles/
|   +-- logs/
|
+-- reports/
|   +-- cng_report.md
|   +-- cnc_report.md
|   +-- final_report.md
|
+-- checksums/
    +-- SHA256SUMS.txt
\end{verbatim}

\subsection{Responsabilités des modules}

\begin{tabularx}{\textwidth}{
    >{\RaggedRight\arraybackslash}p{0.40\textwidth}
    X
}
\toprule
\textbf{Module} & \textbf{Responsabilité} \\
\midrule
\texttt{adaptive\_rbf.py}
& Reprise et adaptation de la classe source. \\

\texttt{neural\_core\_general.py}
& Assemblage des sondes et composeur CNG. \\

\texttt{neural\_core\_calculation.py}
& Sondes dédiées aux opérations. \\

\texttt{synthetic\_data.py}
& Génération et séparation des données. \\

\texttt{operation\_graph.py}
& Définition et exécution des graphes CNC. \\

\texttt{benchmark.py}
& Mesures de temps, mémoire et paramètres. \\

\texttt{audit.py}
& Contrôles d'intégrité et de reproductibilité. \\
\bottomrule
\end{tabularx}

\subsection{Réutilisation du code source}

La classe AdaptiveRBF doit être reprise depuis :

\begin{verbatim}
inputs/reference_v002/src/ex6_common.py
\end{verbatim}

Les éléments à conserver sont :
\begin{itemize}
    \item la structure de l'encodeur ;
    \item la projection latente ;
    \item l'initialisation KMeans ;
    \item le calcul des distances ;
    \item les largeurs apprenables ;
    \item la couche de lecture.
\end{itemize}

Les éléments à adapter sont :
\begin{itemize}
    \item la dimension de sortie ;
    \item les fonctions de perte ;
    \item les générateurs de données ;
    \item les graphes de composition ;
    \item les métriques de calcul numérique.
\end{itemize}

La source originale doit rester inchangée.
Les nouvelles expériences doivent être placées
dans des modules distincts.

% ============================================================
\section{Limites et précautions}
% ============================================================

\subsection{Limites de cette première validation}

Ce protocole ne valide pas encore :

\begin{itemize}
    \item un coordinateur neuronal adaptatif ;
    \item un routage dynamique ;
    \item une sélection conditionnelle de sondes ;
    \item une mémoire temporaire intelligente ;
    \item une modification dynamique du graphe ;
    \item une architecture physique ;
    \item une supériorité générale sur les MLP ;
    \item une supériorité sur le calcul natif.
\end{itemize}

Ces mécanismes devront être étudiés dans des expériences
ultérieures.

\subsection{Limite de la base EX7}

EX7 V003 est une expérience de développement
sur le jeu Wisconsin.

Le rapport final indique notamment :
\begin{itemize}
    \item 1\,341 paramètres persistants pour V003 ;
    \item 97,88 \% d'accuracy ;
    \item 0,06806 de log-loss ;
    \item une latence p50 raw-end-to-end de 0,2902 ms
    au batch 1 ;
    \item une validation externe non établie.
\end{itemize}

Le MLP apparié obtient dans ce rapport :
\begin{itemize}
    \item 98,12 \% d'accuracy ;
    \item 0,05136 de log-loss ;
    \item une latence p50 de 0,0251 ms au batch 1.
\end{itemize}

Ces résultats montrent que la compacité de V003
ne garantit ni une meilleure qualité ni une meilleure
latence.

Ils ne constituent pas une validation du CNG ou du CNC,
qui sont des expériences distinctes.

\subsection{Limite de la spécialisation}

Le fait d'affecter une opération à une sonde ne garantit
pas que cette sonde apprenne une représentation
réutilisable.

Cette propriété doit être testée sur plusieurs graphes.

\subsection{Limite de la composition}

Une composition peut amplifier les erreurs.

Il faut donc mesurer :
\begin{itemize}
    \item l'erreur de chaque opération ;
    \item l'erreur de chaque noeud ;
    \item l'erreur finale ;
    \item la profondeur du graphe ;
    \item le coût d'exécution.
\end{itemize}

% ============================================================
\section{Décision de passage vers l'étape suivante}
% ============================================================

La progression expérimentale est :

\begin{figure}[H]
\centering
\begin{adjustbox}{max width=0.82\textwidth}
\begin{tikzpicture}[node distance=9mm]

\node[block] (p)
{Sonde adaptative};

\node[block, below=of p] (g)
{Coeur neuronal général};

\node[block, below=of g] (c)
{Coeur neuronal de calcul};

\node[block, below=of c] (coord)
{Coordination neuronale apprise};

\node[block, below=of coord] (multi)
{Architecture multi-coeurs};

\draw[flow] (p)--(g);
\draw[flow] (g)--(c);
\draw[flow] (c)--(coord);
\draw[flow] (coord)--(multi);

\end{tikzpicture}
\end{adjustbox}
\caption{Progression expérimentale du projet th\_sonde.}
\end{figure}

Le passage vers l'étape suivante n'est autorisé
que lorsque l'étape précédente a produit des résultats
reproductibles.

Le coordinateur neuronal adaptatif ne doit pas être
introduit avant d'avoir vérifié :
\begin{enumerate}
    \item que les sondes apprennent correctement ;
    \item que plusieurs sondes peuvent être combinées ;
    \item que des sondes spécialisées peuvent être réutilisées ;
    \item que les graphes composés fonctionnent ;
    \item que le coût de composition est quantifié.
\end{enumerate}

% ============================================================
\section{Conclusion}
% ============================================================

Cette expérience constitue le socle expérimental
du projet th\_sonde.

Le \CNG{} vérifie la capacité de plusieurs sondes
adaptatives à apprendre collectivement une fonction.

Le \CNC{} vérifie la capacité de sondes spécialisées
à apprendre des opérations élémentaires et à les
composer dans des graphes de calcul.

Les deux expériences sont complémentaires, mais
ne doivent pas être confondues.

Le CNG répond à la question :

\begin{center}
\textit{Plusieurs sondes peuvent-elles former
un coeur neuronal fonctionnel ?}
\end{center}

Le CNC répond à la question :

\begin{center}
\textit{Des sondes neuronales spécialisées peuvent-elles
réaliser et composer des opérations de calcul ?}
\end{center}

La coordination adaptative, la mémoire et la modification
dynamique des graphes seront étudiées ultérieurement.

\[
\boxed{
\text{Sonde}
\rightarrow
\text{CNG}
\rightarrow
\text{CNC}
\rightarrow
\text{Coordination apprise}
}
\]

La règle fondamentale est de valider chaque niveau
indépendamment avant d'introduire le niveau suivant.

% ============================================================
\appendix
\section{Checklist de reproductibilité}
% ============================================================

\begin{itemize}
    \item[$\square$] Source AdaptiveRBF identifiée et versionnée.
    \item[$\square$] Source originale conservée sans modification.
    \item[$\square$] Configuration expérimentale gelée.
    \item[$\square$] Graines enregistrées.
    \item[$\square$] Domaines mathématiques définis.
    \item[$\square$] Ensembles train, validation et test disjoints.
    \item[$\square$] KMeans limité aux données train.
    \item[$\square$] Nombre de paramètres calculé automatiquement.
    \item[$\square$] Comparateurs entraînés.
    \item[$\square$] Ablations exécutées.
    \item[$\square$] Checkpoints sauvegardés.
    \item[$\square$] Prédictions sauvegardées.
    \item[$\square$] Erreurs intermédiaires CNC sauvegardées.
    \item[$\square$] Latences mesurées.
    \item[$\square$] Résultats par graine conservés.
    \item[$\square$] Résultats négatifs conservés.
    \item[$\square$] Audit final effectué.
    \item[$\square$] Empreintes SHA-256 générées.
\end{itemize}

\end{document}